\section{Thermische Machines en de Tweede Hoofdwet}
\subsection{Van Warmte naar Arbeid}

\begin{definition}
    \textbf{Thermische Arbeidsmachines}: machines die warmte omzetten in arbeid, of omgekeerd.
\end{definition}
\begin{definition}
    \textbf{Warmtepomp}: machine die arbeid gebruikt om warmte van een koud reservoir naar een warm reservoir te transporteren.
\end{definition}

\subsection{Thermische Machines}

De eerste wet zegt ons dat \(|Q_H| - |Q_K| = |W|\).

\begin{definition}
    \textbf{Thermisch Rendement}: \(\eta = \frac{|W|}{|Q_H|} = \frac{|Q_H| - |Q_K|}{|Q_H|} = 1 - \frac{|Q_K|}{|Q_H|}\)
\end{definition}

Twee types machines:
\begin{itemize}
    \item \textbf{Uitwendige verbrandingsmachines}: machines waarbij de verbranding buiten de machine plaatsvindt, zoals stoommachines.
    \item \textbf{Inwendige verbrandingsmachines}: machines waarbij de verbranding binnen de machine plaatsvindt, zoals benzinemotoren.
\end{itemize}

\subsubsection{De Stoommachine}

\begin{definition}
    \textbf{Stoommachine}: een uitwendige verbrandingsmachine die gebruik maakt van de expansie van stoom om arbeid te verrichten.
\end{definition}    

Onze benaderende cyclus is de \textbf{Rankine-cyclus}, bestaande uit volgende zes processen:
\begin{itemize}
    \item 1\(\to\)2: Adiabatische compressie
    \item 2\(\to\)3: Isobare verwarming van water tot kookpunt.
    \item 3\(\to\)4: Isobare en isotherme verdamping.
    \item 4\(\to\)5: Isobare oververhitting van stoom.
    \item 5\(\to\)6: Adiabatische expansie van stoom.
    \item 6\(\to\)1: Isobare en isotherme condensatie van stoom.
\end{itemize}

\subsubsection{De Stirlingmotor}

De geïdealiseerde cyclus van een Stirlingmotor bestaat uit volgende vier processen:
\begin{itemize}
    \item 1\(\to\)2: Isotherme compressie van het gas bij lage temperatuur \(T_K\).
    \item 2\(\to\)3: Isochore warmteopname waarbij het gas wordt opgewarmd van \(T_K\) naar \(T_H\).
    \item 3\(\to\)4: Isotherme expansie van het gas bij hoge temperatuur \(T_H\).
    \item 4\(\to\)1: Isochore warmteafvoer waarbij het gas wordt gekoeld van \(T_H\) naar \(T_K\).
\end{itemize}

De geïdealiseerde cyclus steunt op vijf veronderstellingen:
\begin{itemize}
    \item Het werkmedium is een ideaal gas.
    \item Er zijn geen lekken.
    \item De warmte gaat niet verloren aan de omgeving door de wanden.
    \item De regenerator is een perfecte thermische isolator.
    \item Er is geen wrijving.
\end{itemize}


\subsubsection{De Benzinemotor}

Het viertaktproces van een benzinemotor bestaat uit vier bewegingen. Voor de volledigheid worden zes processen beschreven:
\begin{itemize}
    \item Inlaat (1e beweging): het inlaatventiel opent en het mengsel van lucht en brandstof wordt aangezogen.
    \item Compressie (2e beweging): het mengsel wordt samengeperst, waardoor de temperatuur en druk stijgen.
    \item Ontsteking: het mengsel wordt ontstoken, meestal door een vonk, wat leidt tot een snelle verbranding en expansie van gassen.
    \item Expansie (3e beweging): de hete gassen duwen de zuiger naar beneden, waardoor arbeid wordt verricht.
    \item Eerste uitlaat: het uitlaatventiel opent en de verbrande gassen worden afgevoerd.
    \item Tweede uitlaat (4e beweging): de zuiger beweegt omhoog om de laatste restanten van de verbrande gassen uit te duwen.
\end{itemize}

De bekendste cyclus hiervan is de \textbf{Otto-cyclus}, genaamd achter een vriend van Miel, bestaande uit volgende zes (waarvan vier essentieel) processen:
\begin{itemize}
    \item 1\(\to\)2: Quasistatische adiabatische compressie van het gas.
    \item 2\(\to\)3: Quasistatisch isochore temperatuur en druktoename waarbij het gas wordt verbrand en opgewarmd.
    \item 3\(\to\)4: Quasistatische adiabatische expansie van het gas.
    \item 4\(\to\)1: Quasistatisch isochore temperatuur en drukafname waarbij de verbrande gassen worden afgevoerd.
    \item 5\(\to\)1: Quasistatische isobare inlaat van het lucht-brandstofmengsel.
    \item 1 \(\to\)5: Quasistatische isobare uitlaat van de verbrande gassen bij atmosferische druk.
\end{itemize}

Hier zijn we opnieuw uitgegaan van een aantal veronderstellingen:
\begin{itemize}
    \item Het werkmedium is een ideaal gas.
    \item Alle processen zijn quasistatisch.
    \item Er is geen wrijving.
\end{itemize}

Processen 1 \(\to\) 5 en 5 \(\to \) 1 compenseren elkaar en wordt niet meer in rekening gebracht bij de berekening van het rendement.
Er zijn namelijk nog twee processen waar er warmte wordt uitgewisseld, namelijk 2 \(\to\) 3 en 4 \(\to\) 1. 
We kunnen schrijven:
\[Q_H = C_V (T_3-T_2) = |Q_H|\]
en 
\[Q_K = C_V (T_1-T_4) = -|Q_K|\]
Hieruit volgt dat het rendement van de Otto-cyclus gelijk is aan: 
\[\eta = 1 - \frac{|Q_K|}{|Q_H|} = 1 - \frac{T_4-T_1}{T_3-T_2}\]

Maar we willen een uitdrukking in functie van het volume in plaats van de temperatuur, want het volume is makkelijker te beheersen.
We kunnen hiervoor gebruik maken van de adiabatische relaties van de processen 1\(\to\)2 en 3\(\to\)4:
\[T_1 V_1^{\gamma - 1} = T_2 V_2^{\gamma - 1} \quad \text{en} \quad T_4 V_1^{\gamma - 1} = T_3 V_2^{\gamma - 1}\]
Hieruit volgt dat:
\[(T_4-T_1)V_1^{\gamma - 1} = (T_3-T_2)V_2^{\gamma - 1}\]
\[\implies \frac{T_4-T_1}{T_3-T_2} = \left(\frac{V_2}{V_1}\right)^{\gamma - 1}\]
Met \(r = \frac{V_1}{V_2}\) als compressieverhouding, kunnen we het rendement van de Otto-cyclus schrijven als:
\[\eta = 1 - \frac{|Q_K|}{|Q_H|} = 1 - \left(\frac{V_2}{V_1}\right)^{\gamma - 1} = 1 - \frac{1}{r^{\gamma - 1}}\]

\subsection{De Dieselmotor}

Er zijn opnieuw zes processen waarvan er vier nodig zijn voor de cyclus waar we geïnteresseerd in zijn:
Waar we de volgende effecten genegeren: chemische reacties, warmteverlies, wrijving, niet-ideale gassen, niet-quasistatische processen, enzovoort.   

Onze cyclus is nu de \textbf{lucht-gestandardiseerde dieselcyclus}, bestaande uit volgende zes (waarvan vier essentieel) processen:
\begin{itemize}
    \item 1\(\to\)2: Quasistatische adiabatische compressie van het gas.
    \item 2\(\to\)3: Quasistatisch isobare temperatuur en volume toename waarbij het gas wordt verbrand en opgewarmd.
    \item 3\(\to\)4: Quasistatische adiabatische expansie van het gas.
    \item 4\(\to\)1: Quasistatisch isochore temperatuur en drukafname waarbij de verbrande gassen worden afgevoerd.
    \item 5\(\to\)1: Quasistatische isobare inlaat van lucht.
    \item 1 \(\to\)5: Quasistatische isobare uitlaat van de verbrande gassen bij atmosferische druk.
\end{itemize}

Het grote verschil met de Otto-cyclus is dat het verbrandingsproces (2\(\to\)3) nu isobaar is in plaats van isochore. In deze stap zal 
\[Q_H = C_P(T_3 -T_2) = |Q_H|\]
waardoor het rendement van de dieselcyclus gelijk is aan:
\[\eta = 1 - \frac{|Q_K|}{|Q_H|} = 1 - \frac{C_V(T_4-T_1)}{C_P(T_3 - T_2)} = 1 - \frac{T_4-T_1}{\gamma (T_3 - T_2)}\]
We kunnen ook hier een uitdrukking in functie van het volume bekomen door gebruik te maken van de adiabatische relaties van de processen 1\(\to\)2 en 3\(\to\)4:
\[T_1 V_1^{\gamma - 1} = T_2 V_2^{\gamma - 1} \quad \text{en} \quad T_4 V_1^{\gamma - 1} = T_3 V_3^{\gamma - 1}\]
Hieruit volgt dat:
\[(T_4-T_1)V_1^{\gamma - 1} = T_3V_3^{\gamma - 1}-T_2V_2^{\gamma - 1}\]
\[\implies T_4-T_1 = \frac{T_3V_3^{\gamma - 1}-T_2V_2^{\gamma - 1}}{V_1^{\gamma - 1}}\]
Omdat we een isobaar hebben van 2\(\to\)3, kunnen we schrijven dat \(T_3/T_2 = V_3/V_2 \implies T_2 = T_3 \frac{V_2}{V_3}\). Hieruit volgt dat:
\[\eta = 1 - \frac{\frac{V_3^{\gamma-1}}{V_1^{\gamma-1} }- \frac{V_2^{\gamma-1}}{V_3V_1^{\gamma-1}}}{1 - \frac{V_2}{V_3}} 
= 1 - \frac{\left(\frac{V_3}{V_1}\right)^\gamma-\left(\frac{V_2}{V_1}\right)^{\gamma}}{\frac{V_3}{V_1}- \frac{V_2}{V_1}}\]
Als we de \textbf{expansieverhouding}:
\[r_E = \frac{V_1}{V_3} \]
en de \textbf{compressieverhouding}:
\[r_C = \frac{V_1}{V_2}\]
definiëren, kunnen we het rendement van de dieselcyclus schrijven als:
\[
\eta = 1 - \frac{1}{\gamma} \frac{(\frac{1}{r_E})^{\gamma} - (\frac{1}{r_C})^{\gamma}}{\frac{1}{r_E} - \frac{1}{r_C}} 
= 1 - \frac{1}{\gamma} \frac{r_E^{-\gamma} - r_C^{-\gamma}}{r_E^{-1} - r_C^{-1}}
\]

\subsection{Tweede Wet: Kelvin-Planck Formulering}

\begin{definition}
    \textbf{warme bron}: een systeem dat in staat is om warmte af te geven zonder zelf van toestand te veranderen.
\end{definition}
\begin{definition}
    \textbf{koude bron}: een systeem dat in staat is om warmte op te nemen zonder zelf van toestand te veranderen.
\end{definition}
\begin{postulate}
    Het is onmogelijk om een cyclus te hebben waarvan het enige resultaat is dat er warmte wordt onttrokken aan een warme bron en volledig wordt omgezet in arbeid. 
\end{postulate}
\begin{postulate}
    \textbf{Kelvin formulering}: "Het is onmogelijk door tussenkomst van niet-levende materie, mechanische effecten te bekomen uit een deel van deze materie, door ze af te koelen beneden de temperatuur van het koudste object in de omgeving."
\end{postulate}
\begin{postulate}
    \textbf{Planck formulering}: "Het is onmogelijk een cyclus te construeren waarvan het enige resultaat is dat er warmte wordt onttrokken aan een warme bron (warmtereservoir) en volledig wordt omgezet in arbeid."
\end{postulate}

\begin{postulate}
    \textbf{Tweede hoofdwet van de thermodynamica}: Er is geen enkel proces mogelijk waarvan het enige resultaat is dat er warmte wordt onttrokken aan een warme bron en volledig wordt omgezet in arbeid.
\end{postulate}

Een machine die de eerste hoofdwet overtreedt is een \textbf{perpetuum mobile van de eerste soort}. Een machine die de tweede hoofdwet overtreedt is een \textbf{perpetuum mobile van de tweede soort}.

\subsection{Warmtepompen}

\begin{definition}
    \textbf{Prestatiecoëfficiënt}: \(\omega= \frac{|Q_K|}{|W|} = \frac{|Q_K|}{|Q_H| - |Q_K|} = \frac{|Q_H|}{|W|} - 1\)
\end{definition}

\subsubsection{De Stirlingwarmtepomp}

De cyclus van de Stirlingmotor is reversibel, dus kunnen we deze ook gebruiken als warmtepomp.
De vier processen zijn nu:
\begin{itemize}
    \item 1\(\to\)2: Isotherme compressie van het gas bij hoge temperatuur \(T_H\).
    \item 2\(\to\)3: Isochore warmteafvoer waarbij het gas wordt gekoeld van \(T_H\) naar \(T_K\).
    \item 3\(\to\)4: Isotherme expansie van het gas bij lage temperatuur \(T_K\).
    \item 4\(\to\)1: Isochore warmteopname waarbij het gas wordt opgewarmd van \(T_K\) naar \(T_H\).
\end{itemize}

\subsubsection{De Frigo (Koelkast)}
De meeste koelmachines steunen op het \textbf{Joule-Kelvin-smoorproces}, waar men op de eigenschap van gesatureerde vloeistoffen steunt.
Het proces van de geïdealiseerde koelmachine, dus wanneer wrijving, warmteverlies, turbulentie, enzovoort worden genegeerd, bestaat uit volgende stappen:
\begin{itemize}
    \item 1\(\to\)2: Smoorproces met druk en temperatuurafname van het vloeibare koelmiddel. (Kan hier niet met thermodynamische coördinaten beschreven worden.) 
    \item 2\(\to\)3: Isotherme en isobare verdamping waarbij een hoeveelheid warmte \(Q_K\) onttrokken wordt.
    \item 3\(\to\)4: Adiabatische compressie waarbij het gasvormige koelmiddel wordt samengeperst tot een hoge druk en temperatuur.
    \item 4\(\to\)1: Isobare afkoeling waarbij het koelmiddel condenseert en warmte \(Q_H\) afgeeft aan de omgeving.
\end{itemize}

\subsection{De formulering van  Clausius voor de Tweede Hoofdwet}

\begin{postulate}
    \textbf{Clausius formulering}: "Geen enkel proces is mogelijk waarvan het enige resultaat is dat er warmte van een koud reservoir naar een warm reservoir wordt getransporteerd."
\end{postulate} 

\subsection{Gelijkwaardigheid Kelvin-Planck (K) en Clausiusformulering (C) van de Tweede Hoofdwet}

\begin{derivation}
    \textbf{Bewijs van de gelijkwaardigheid van de Kelvin-Planck en Clausius formulering van de tweede hoofdwet.}
\begin{itemize}
    \item \textbf{K \(\implies\) C}: 
    
    Stel dat een arbeidsmachine niet voldoet aan de Kelvin-Planck formulering, dan moeten we bewijzen dat ze ook niet voldoet aan de Clausius formulering.

    Stel dat er een proces is waarbij warmte van een koud reservoir naar een warm reservoir wordt getransporteerd zonder arbeid te verrichten. 
    Er wordt dus alleen een hoeveelheid warmte \(Q_1\) onttrokken van de warme bron en een arbeid \(W\) op de omgeving verricht.

        \(|Q_1| =  |W|\) 
    
    Deze arbeid wordt nu gebruikt om een warmtepomp die een hoeveelheid warmte \(Q_2\) onttrekt van een koud reservoir en een hoeveelheid warmte afgeeft aan een warm reservoir te laten werken.
    We kunnen volgens de eerste wet schrijven dat:
    \[|Q'_H | = |Q_2| + |W| = |Q_1| + |Q_2|\]
    Er is dus een netto warmteoverdracht van \(Q_2\) van het koude reservoir naar het warme reservoir, wat in strijd is met de Clausius formulering van de tweede hoofdwet.
    
    \item \textbf{C \(\implies\) K}: 
    
    Stel dat een machine niet voldoet aan de Clausius formulering, dan moeten we bewijzen dat ze ook niet voldoet aan de Kelvin-Planck formulering.

    Stel dat er een machine is die de Kelvin-Planck formulering overtreedt, dus een machine die warmte volledig omzet in arbeid. 
    Veronderstel dat de arbeidsmachine een hoeveelheid warmte \(Q_1\) onttrekt van een warme bron en een hoeveelheid arbeid \(W\) verricht, waarbij ze een hoeveelheid warmte \(|Q_2|\) afstaat aan de koude bron.
    Dan zegt de eerste hoofdwet dat:
    \[|Q_1| = |W| + |Q_2| \implies |W| = |Q_1| - |Q_2|\]
    Het netto resultaat van deze machine is dat er een hoeveelheid warmte \(|Q_1|-|Q_2|\) onttrokken wordt van een warme bron en volledig wordt omgezet in arbeid, wat in strijd is met de Kelvin-Planck formulering van de tweede hoofdwet.
\end{itemize}
\qed
\end{derivation}